\section*{Chapter \ref{ch:ll:testing-and-deploying-low-latency-systems} --- Testing and Deploying Low-Latency Systems}

\begin{solution}{exo:ll:testing-and-deploying-low-latency-systems:1}
No order ever has more filled than it asked for; the position never exceeds the limit counted in the worst case (or: the
position equals the sum of the fills). That the gateway's view agrees with the venue's records needs an independent record of
the venue's side, the \omterm{def:ll:order-gateway-and-order-management:dropcopy}{drop copy}.
\end{solution}

\begin{solution}{exo:ll:testing-and-deploying-low-latency-systems:2}
So that a read one byte past the input is a read past an allocation, which AddressSanitizer reports. Inside a larger buffer
(a string's capacity, a reused array) the same bad read would return garbage silently.
\end{solution}

\begin{solution}{exo:ll:testing-and-deploying-low-latency-systems:3}
Anything that changes with time (another job starting, the processor warming up, a frequency change) then affects one build and
not the other, and is measured as a difference between them. Interleaving spreads drifts over both.
\end{solution}

\begin{solution}{exo:ll:testing-and-deploying-low-latency-systems:4}
$2\,(1.645 + 1.282)^2\,(1.2533 \times 0.14 / r)^2$: about 53 runs per side for $r = 10\%$, about 840 for $r = 2.5\%$. The
number grows as the inverse square of the rise.
\end{solution}

\begin{solution}{exo:ll:testing-and-deploying-low-latency-systems:5}
$20 \times 21 = 420$ gates; at a level of 5\% with no real change, up to 21 false alarms. The tolerance flags only rises larger
than 2\%, which removes many of them: the measured rate on this laptop was 3 to 9\%, some 13 to 38 a month; a larger tolerance or
a lower level would cut them, at the cost of power. Each false alarm costs a rerun, so the level and the tolerance are set with the team's patience in mind.
\end{solution}

\begin{solution}{exo:ll:testing-and-deploying-low-latency-systems:6}
The checksum is the sum of the bytes modulo 256. Turning ``1'' into ``0'' lowers the sum by one, turning ``2'' into ``3'' raises
it by one: the sum is unchanged. It catches most single-byte corruptions and truncations; it does not catch compensating
changes, reorderings of bytes, or anything deliberate. The protection against corruption on the wire is TCP's; the parser's
checks are about structure.
\end{solution}

\begin{solution}{exo:ll:testing-and-deploying-low-latency-systems:7}
A venue configured with a message rate (the simulator's throttle), a burst of orders sent at the same time stamp, and
expectations for a sequence of acceptances followed by rejections with reason \texttt{T}; the step depends on time, so the
runner must control the time stamps and the engine's configuration, which the other steps did not need.
\end{solution}

\begin{solution}{exo:ll:testing-and-deploying-low-latency-systems:8}
One run of each, a week apart: the machine, its load, its temperature and the rest of the system differ between the two, and
on this laptop a single run's 99th percentile moves by far more than 5\%. The gate would fail at random and pass real
regressions. Run both builds now, interleaved, many times, and compare them with a test.
\end{solution}

\begin{solution}{pb:ll:testing-and-deploying-low-latency-systems:1}
\begin{enumerate}
\item About \qty{208}{\nano\second}, with a robust coefficient of variation of about 14\%.
\item About 18\% (\qty{245}{\nano\second} against 208): the planted wait of \qty{11.4}{\nano\second} cost more than itself,
because the time stamp reads of the wait take time too.
\item It is set by the slowest 1\% of the run's events, those hit by interruptions and the system's own work, so it depends on
what the rest of the machine did during that run.
\item Its difference varies by about 20\% from pair to pair, four times the rise to detect.
\item Medians ignore the runs spoiled by a burst of other work, which would pull a mean.
\item That, without a real difference, the labels A and B could be swapped (the interleaving makes this plausible); not that
the runs are normal, or that A and B have the same spread.
\item 3 to 9\% over a hundred gates, around the test's level of 5\% (a hundred gates estimate a rate of 5\% to within a few
points); the tolerance removes the significant rises below 2\%.
\item 40.
\item 57\%, 73\% and 84\%.
\item About 200 runs per side, where the rate reaches 91\%.
\item About 220, close: the approximation assumes normal runs and uses the median's efficiency for them; the measured runs
are concentrated with a few far outliers, for which the median does a little better than that.
\item About $2 \times 200 \times \qty{30}{\milli\second} \approx \qty{12}{\second}$.
\item \textbf{Named result.} On this laptop, whose runs' 99th percentiles vary by about 14\% even when it is otherwise idle,
the gate needs about 200 interleaved runs per side (about 12 seconds of benchmark) to detect a 5\% rise with 90\% power, at a
false-alarm rate near the test's 5\% (the normal approximation says about 220); an 18\% regression needs 40.
\item Less variation from run to run; the runs needed fall as the square of the spread, to about 27 at a spread of 5\%.
\item With enough runs, any change becomes significant; the tolerance says which changes matter.
\item Far fewer runs, since medians are stable; but it misses what the 99th percentile catches: a rare slow path, such as an
occasional allocation or a log line on one branch.
\item Run it on candidates for merge and nightly, on a dedicated machine, with the number of runs set by the machine's
measured noise; keep the functional tests on every change.
\item Changes in the network, the kernel, the interaction with other processes, the hardware, and data that differ from the
benchmark's day.
\item To recompute the verdict and the power analysis with other statistics or thresholds, and to compare machines.
\item Enough interleaved runs for the machine's measured noise, judged by a test that assumes little, and a tolerance.
\end{enumerate}
\end{solution}

\begin{solution}{iq:ll:testing-and-deploying-low-latency-systems:1}
A test that generates random inputs and checks a property for all of them. For an order book: after any sequence of adds,
modifications and cancels, the best bid is below the best ask, the quantity at each level equals the sum of its orders, and
removing every order leaves an empty book.
\iqlookfor{a property, a generator, and shrinking of failures.}
\end{solution}

\begin{solution}{iq:ll:testing-and-deploying-low-latency-systems:2}
Mutate recorded packets (bytes, lengths, sequence numbers, truncations), feed them in exact-size buffers under sanitisers, and
check invariants besides crashes: every accepted message within its buffer, sequence numbers consistent, the book valid after
each message, and agreement with a reference parser.
\iqlookfor{sanitisers plus invariants and a differential reference.}
\end{solution}

\begin{solution}{iq:ll:testing-and-deploying-low-latency-systems:3}
Ask how it was measured: interleaved runs of both builds on the same machine, enough of them for its noise, a test with a
confidence interval. Rerun with more runs if needed; 3\% may be within the noise, or below the tolerance even if real.
\iqlookfor{interleaving, noise, a test and a tolerance.}
\end{solution}

\begin{solution}{iq:ll:testing-and-deploying-low-latency-systems:4}
That the client's messages and its handling of the venue's responses follow the protocol, step by step, on the venue's test
system. Not the firm's strategy, its risk controls beyond what the script exercises, its latency, or its behaviour under load
and failures in production.
\iqlookfor{protocol conformance, not correctness or performance.}
\end{solution}

\begin{solution}{iq:ll:testing-and-deploying-low-latency-systems:5}
Preflight the host; run the new build in shadow on live inputs and compare its outputs; start it on a small share (a canary),
with tight limits, then widen; roll back by redeploying the previous build, a procedure prepared and tested in advance, with the
journal and positions carried over.
\iqlookfor{shadow, canary, prepared rollback.}
\end{solution}

\begin{solution}{iq:ll:testing-and-deploying-low-latency-systems:6}
On every change: unit, property and \omterm{def:ll:testing-and-deploying-low-latency-systems:tests}{differential tests}, and fuzzing with sanitisers. On candidates: simulator scenarios and
replays of recorded days, certification scripts, and the performance gate on a dedicated machine. On release: host preflight,
shadow running, canary, staged rollout, prepared rollback. Each stage catches what the one before cannot see.
\iqlookfor{the pyramid, and what each layer catches.}
\end{solution}
